% Appendix 8: Vector-Valued Motion in Three Dimensions
% VIC Specialist Mathematics — Units 3–4;
% QLD Specialist Mathematics — Unit 3 (vector calculus)
\begingroup
\small

\noindent\fbox{%
\begin{minipage}{0.97\textwidth}
\textit{\textbf{Enrichment only.}
Three-dimensional vector calculus is part of
\textbf{VIC Specialist Mathematics (Units~3--4)} and
\textbf{QLD Specialist Mathematics (Unit~3)}.
NSW Extension~1 already covers vector functions for motion in a plane;
this appendix extends those ideas by adding a third component.}
\end{minipage}}

\vspace{0.8em}

% ---- Key concepts ----
\noindent\textbf{Key concepts.}

The position vector of a moving particle is
\[
  \boxed{
  \mathbf{r}(t)
  = x(t)\,\mathbf{i} + y(t)\,\mathbf{j} + z(t)\,\mathbf{k}
  }.
\]
The \textbf{velocity} vector is the derivative of position:
\[
  \boxed{\mathbf{v}(t) = \mathbf{r}'(t)}.
\]
The \textbf{acceleration} vector is the derivative of velocity:
\[
  \boxed{\mathbf{a}(t) = \mathbf{r}''(t)}.
\]
The particle's \textbf{speed} is the magnitude of its velocity:
\[
  \boxed{
  |\mathbf{v}(t)|
  = \sqrt{x'(t)^{2} + y'(t)^{2} + z'(t)^{2}}
  }.
\]
The same differentiation rules used in two dimensions apply
component by component in three dimensions.

% ---- Worked example ----
\vspace{0.8em}
\noindent\textbf{Worked example.}

\noindent A particle moves through space with position
\[
  \mathbf{r}(t) = t\,\mathbf{i} + t^{2}\,\mathbf{j}
                + t^{3}\,\mathbf{k},
  \qquad t \ge 0.
\]
Find: (a) its velocity and acceleration;\quad
      (b) its speed at $t = 1$;\quad
      (c) the equation of the tangent line to its path at $t = 1$.

\vspace{0.3em}
\noindent\textbf{(a) Velocity and acceleration.}

Differentiating each component,
\[
  \boxed{
  \mathbf{v}(t) = \mathbf{i} + 2t\,\mathbf{j} + 3t^{2}\,\mathbf{k}
  }.
\]
Differentiating again,
\[
  \boxed{
  \mathbf{a}(t) = 2\,\mathbf{j} + 6t\,\mathbf{k}
  }.
\]

\vspace{0.3em}
\noindent\textbf{(b) Speed at $\boldsymbol{t = 1}$.}

The velocity at $t = 1$ is
\[
  \mathbf{v}(1) = (1, 2, 3).
\]
Therefore,
\[
  |\mathbf{v}(1)|
  = \sqrt{1^{2} + 2^{2} + 3^{2}}
  = \boxed{\sqrt{14}}.
\]

\vspace{0.3em}
\noindent\textbf{(c) Tangent line.}

At $t = 1$, the particle is at
\[
  \mathbf{r}(1) = (1, 1, 1).
\]
The velocity vector $(1, 2, 3)$ gives the direction of the tangent.
Hence the tangent line is
\[
  \boxed{
  \mathbf{R} = (1, 1, 1) + \lambda\,(1, 2, 3),
  \qquad \lambda \in \mathbb{R}
  }.
\]
This connects the familiar NSW vector equation of a line with
calculus: the velocity gives the tangent direction to the particle's
path.

% ---- Key takeaways ----
\vspace{0.8em}
\noindent\textbf{Key takeaways.}
\begin{itemize}[leftmargin=*,itemsep=2pt]
  \item Velocity is a vector; speed is its magnitude.
        They are not interchangeable.
  \item Differentiation of vector functions works component by
        component.
  \item The velocity vector gives the instantaneous tangent direction
        to a particle's path, even in three-dimensional space.
\end{itemize}

\vspace{0.4em}
\noindent\textit{Beyond school --- University connection.}
Three-dimensional vector calculus introduces space curves, tangent
vectors, arc length and curvature
(see also Problem~2.20 in Part~1 of this booklet).
These ideas are used in robotics, aerospace engineering, mechanics
and computer animation.

% ================================================================
%  Final Thoughts: Three Operations, Three Geometric Ideas
% ================================================================
\vspace{1.2em}
\noindent\rule{\textwidth}{0.5pt}
\vspace{0.4em}

\noindent\textbf{\large Final Thoughts: Three Operations,
Three Geometric Ideas}

\vspace{0.4em}
\noindent
The dot product, cross product and linear independence describe
different aspects of vector geometry.

\vspace{0.5em}
\begin{center}
\renewcommand{\arraystretch}{1.4}
\begin{tabular}{|p{0.30\textwidth}|p{0.58\textwidth}|}
\hline
\textbf{Concept} & \textbf{What it tells us} \\
\hline
Dot product $\mathbf{a} \cdot \mathbf{b}$
  & Angles, projections and perpendicularity \\
\hline
Cross product $\mathbf{a} \times \mathbf{b}$
  & A perpendicular vector and an area \\
\hline
Linear independence
  & Whether vectors provide genuinely different directions \\
\hline
\end{tabular}
\end{center}

\vspace{0.5em}
\noindent
Together, these concepts explain how vectors describe directions,
lines, planes and motion.

\vspace{0.4em}
\noindent
The underlying message is simple:

\begin{center}
\itshape
Vectors are more than arrows.
They provide a common mathematical language for geometry, motion,
linear algebra and many problems encountered in university STEM
courses.
\end{center}

\endgroup
